% GENERATED FILE. Edit canonical lessons, then run npm run generate:books.
% canonical-source-hash: fnv1a64:d48189bc66074f2c
% canonical-lessons: TE-C460-asaukaryam, TE-C460-ibbandi, TE-C460-vakabu-ceyu, TE-C460-vapasu, TE-C460-saricusuko

\chapter{Complaints at the Desk}
\label{ch:te-complaints-at-the-desk}
\hlchaptermodality{460}

\begin{chapteropening}
\textbf{By the end of this chapter:} \emph{I can say inconvenience, trouble, to inquire, a refund and to double-check.}

Complete the last lesson of chapter 460: I can say inconvenience, trouble, to inquire, a refund and to double-check.
\end{chapteropening}

\section[\texorpdfstring{\te{అసౌకర్యం} (asaukaryaṁ)}{అసౌకర్యం (asaukaryaṁ)}]{\te{అసౌకర్యం} (asaukaryaṁ) — inconvenience}

\label{lesson:TE-C460-asaukaryam}



\begin{quote}
Before the new one: say the Telugu for the web, then the Telugu for a phone charger.
\end{quote}

\subsection*{You'll want to know: \te{అసౌకర్యం}}
\textbf{\te{అసౌకర్యం}} — \emph{asaukaryaṁ} — \textquotedblleft{}inconvenience\textquotedblright{}.

\begin{quote}
\te{ఆలస్యం} \te{వల్ల} \te{కలిగిన} \te{అసౌకర్యానికి} \te{క్షమించండి}.
\end{quote}

\begin{quote}
\emph{ālasyaṁ valla kaligina asaukaryāniki kṣamiñcaṇḍi.}
\end{quote}

\begin{quote}
We are sorry for the inconvenience caused by the delay.
\end{quote}

\te{అ}-, not, + \te{సౌకర్యం}, comfort. Station announcements use this sentence.

\begin{cousinweb}
\textbf{\te{అసౌకర్యం}} (\emph{asaukaryaṁ}) is built from Sanskrit \emph{a-}, \textquotedblleft{}not\textquotedblright{}, and \emph{saukarya}, \textquotedblleft{}ease, practicability\textquotedblright{} (Monier-Williams).
\end{cousinweb}

\begin{grammarlens}[title={What you've built}]
One more word for getting around.
\end{grammarlens}

\subsection*{Your turn}
\begin{itemize}
\raggedright
  \item \emph{Say it:} \emph{asaukaryaṁ}
  \item \emph{Say it:} \emph{asaukaryaṁ}, once more
  \item \emph{Say it:} \emph{ālasyaṁ valla kaligina asaukaryāniki kṣamiñcaṇḍi.} — We are sorry for the inconvenience caused by the delay.
  \item \emph{Recall:} say the Telugu for the web, then the Telugu for a phone charger, then say \emph{asaukaryaṁ} again
\end{itemize}

\begin{tcolorbox}[breakable,colback=teal!4,colframe=teal!35!black,title={Before you move on}]
What does \te{అసౌకర్యం} mean? (\textquotedblleft{}Inconvenience\textquotedblright{}.) Say it once more.
\end{tcolorbox}
\section[\texorpdfstring{\te{ఇబ్బంది} (ibbandi)}{ఇబ్బంది (ibbandi)}]{\te{ఇబ్బంది} (ibbandi) — trouble; difficulty}

\label{lesson:TE-C460-ibbandi}



\begin{quote}
Before the new one: say the Telugu for a phone charger, then the Telugu for inconvenience.
\end{quote}

\subsection*{You'll want to know: \te{ఇబ్బంది}}
\textbf{\te{ఇబ్బంది}} — \emph{ibbandi} — \textquotedblleft{}trouble; difficulty\textquotedblright{}.

\begin{quote}
\te{కొత్త} \te{నగరంలో} \te{భాష} \te{వల్ల} \te{ఇబ్బంది} \te{పడ్డాం}.
\end{quote}

\begin{quote}
\emph{kotta nagaraṁlō bhāṣa valla ibbandi paḍḍāṁ.}
\end{quote}

\begin{quote}
In the new city we had trouble because of the language.
\end{quote}

\te{ఇబ్బంది} \te{పడు}, to have trouble; \te{ఇబ్బంది} \te{పెట్టు}, to cause someone trouble.

\begin{cousinweb}
\textbf{\te{ఇబ్బంది}} (\emph{ibbandi}) is native Dravidian: DEDR 524 lists Telugu \emph{ibbandi}, \textquotedblleft{}difficulty, inconvenience\textquotedblright{}, beside \emph{iṟuku}, \textquotedblleft{}to be squeezed, to be stuck in\textquotedblright{}, and \emph{ikkaṭṭu}, \textquotedblleft{}hardship\textquotedblright{}. A tight spot became a difficulty.
\end{cousinweb}

\begin{grammarlens}[title={What you've built}]
One more word for getting around.
\end{grammarlens}

\subsection*{Your turn}
\begin{itemize}
\raggedright
  \item \emph{Say it:} \emph{ibbandi}
  \item \emph{Say it:} \emph{ibbandi}, once more
  \item \emph{Say it:} \emph{kotta nagaraṁlō bhāṣa valla ibbandi paḍḍāṁ.} — In the new city we had trouble because of the language.
  \item \emph{Recall:} say the Telugu for a phone charger, then the Telugu for inconvenience, then say \emph{ibbandi} again
\end{itemize}

\begin{tcolorbox}[breakable,colback=teal!4,colframe=teal!35!black,title={Before you move on}]
What does \te{ఇబ్బంది} mean? (\textquotedblleft{}Trouble; difficulty\textquotedblright{}.) Say it once more.
\end{tcolorbox}
\section[\texorpdfstring{\te{వాకబు} \te{చేయు} (vākabu cēyu)}{వాకబు చేయు (vākabu cēyu)}]{\te{వాకబు} \te{చేయు} (vākabu cēyu) — to inquire; to make enquiries}

\label{lesson:TE-C460-vakabu-ceyu}



\begin{quote}
Before the new one: say the Telugu for inconvenience, then the Telugu for trouble.
\end{quote}

\subsection*{You'll want to know: \te{వాకబు} \te{చేయు}}
\textbf{\te{వాకబు} \te{చేయు}} — \emph{vākabu cēyu} — \textquotedblleft{}to inquire; to make enquiries\textquotedblright{}.

\begin{quote}
\te{రైలు} \te{ఎప్పుడు} \te{వస్తుందో} \te{స్టేషన్లో} \te{వాకబు} \te{చేశాను}.
\end{quote}

\begin{quote}
\emph{railu eppuḍu vastundō sṭēṣanlō vākabu cēśānu.}
\end{quote}

\begin{quote}
I inquired at the station when the train would come.
\end{quote}

Past (I …): \te{వాకబు} \te{చేశాను}, \emph{vākabu cēśānu}.

A careful word for asking after something: a train, a person, a parcel.

\begin{cousinweb}
\textbf{\te{వాకబు} \te{చేయు}} (\emph{vākabu cēyu}) is a Telugu phrase: \te{వాకబు}, \textquotedblleft{}an enquiry\textquotedblright{}, whose origin is not settled, and the native \te{చేయు}, \textquotedblleft{}to do\textquotedblright{} (DEDR 1957).
\end{cousinweb}

\begin{grammarlens}[title={What you've built}]
One more word for getting around.
\end{grammarlens}

\subsection*{Your turn}
\begin{itemize}
\raggedright
  \item \emph{Say it:} \emph{vākabu cēyu}
  \item \emph{Say it:} \emph{vākabu cēyu}, once more
  \item \emph{Say it:} \emph{railu eppuḍu vastundō sṭēṣanlō vākabu cēśānu.} — I inquired at the station when the train would come.
  \item \emph{Recall:} say the Telugu for inconvenience, then the Telugu for trouble, then say \emph{vākabu cēyu} again
\end{itemize}

\begin{tcolorbox}[breakable,colback=teal!4,colframe=teal!35!black,title={Before you move on}]
What does \te{వాకబు} \te{చేయు} mean? (\textquotedblleft{}To inquire; to make enquiries\textquotedblright{}.) Say it once more.
\end{tcolorbox}
\section[\texorpdfstring{\te{వాపసు} (vāpasu)}{వాపసు (vāpasu)}]{\te{వాపసు} (vāpasu) — a refund; a return}

\label{lesson:TE-C460-vapasu}



\begin{quote}
Before the new one: say the Telugu for trouble, then the Telugu for to inquire.
\end{quote}

\subsection*{You'll want to know: \te{వాపసు}}
\textbf{\te{వాపసు}} — \emph{vāpasu} — \textquotedblleft{}a refund; a return\textquotedblright{}.

\begin{quote}
\te{రద్దు} \te{చేసిన} \te{టిక్కెట్టుకు} \te{సగం} \te{డబ్బు} \te{వాపసు} \te{వచ్చింది}.
\end{quote}

\begin{quote}
\emph{raddu cēsina ṭikkeṭṭuku sagaṁ ḍabbu vāpasu vaccindi.}
\end{quote}

\begin{quote}
For the cancelled ticket half the money came back as a refund.
\end{quote}

\te{వాపసు} \te{ఇచ్చు}, to give back; \te{వాపసు} \te{వచ్చు}, to come back.

\begin{cousinweb}
\textbf{\te{వాపసు}} (\emph{vāpasu}) is a loan from Persian \emph{vāpas}, \textquotedblleft{}back, returned\textquotedblright{}, which reached Telugu through Hindustani.
\end{cousinweb}

\begin{grammarlens}[title={What you've built}]
One more word for getting around.
\end{grammarlens}

\subsection*{Your turn}
\begin{itemize}
\raggedright
  \item \emph{Say it:} \emph{vāpasu}
  \item \emph{Say it:} \emph{vāpasu}, once more
  \item \emph{Say it:} \emph{raddu cēsina ṭikkeṭṭuku sagaṁ ḍabbu vāpasu vaccindi.} — For the cancelled ticket half the money came back as a refund.
  \item \emph{Recall:} say the Telugu for trouble, then the Telugu for to inquire, then say \emph{vāpasu} again
\end{itemize}

\begin{tcolorbox}[breakable,colback=teal!4,colframe=teal!35!black,title={Before you move on}]
What does \te{వాపసు} mean? (\textquotedblleft{}A refund; a return\textquotedblright{}.) Say it once more.
\end{tcolorbox}
\section[\texorpdfstring{\te{సరిచూసుకో} (saricūsukō)}{సరిచూసుకో (saricūsukō)}]{\te{సరిచూసుకో} (saricūsukō) — to double-check}

\label{lesson:TE-C460-saricusuko}



\begin{quote}
Before the new one: say the Telugu for to inquire, then the Telugu for a refund.
\end{quote}

\subsection*{You'll want to know: \te{సరిచూసుకో}}
\textbf{\te{సరిచూసుకో}} — \emph{saricūsukō} — \textquotedblleft{}to double-check\textquotedblright{}.

\begin{quote}
\te{బయలుదేరే} \te{ముందు} \te{టిక్కెట్లు} \te{సరిచూసుకున్నాను}.
\end{quote}

\begin{quote}
\emph{bayaludērē mundu ṭikkeṭlu saricūsukunnānu.}
\end{quote}

\begin{quote}
Before setting out I double-checked the tickets.
\end{quote}

Past (I …): \te{సరిచూసుకున్నాను}, \emph{saricūsukunnānu}.

\te{సరి}, right, + \te{చూసు}, to see, + -\te{కో}: to see for oneself that it is right.

\begin{cousinweb}
\textbf{\te{సరిచూసుకో}} (\emph{saricūsukō}) is native Dravidian: \te{సరి}, \textquotedblleft{}right, equal\textquotedblright{}, \te{చూచు}, \textquotedblleft{}to see\textquotedblright{} (DEDR 2735), and -\te{కో}, \textquotedblleft{}for oneself\textquotedblright{} (DEDR 2151).
\end{cousinweb}

\begin{grammarlens}[title={What you've built}]
One more word for getting around.
\end{grammarlens}

\subsection*{Your turn}
\begin{itemize}
\raggedright
  \item \emph{Say it:} \emph{saricūsukō}
  \item \emph{Say it:} \emph{saricūsukō}, once more
  \item \emph{Say it:} \emph{bayaludērē mundu ṭikkeṭlu saricūsukunnānu.} — Before setting out I double-checked the tickets.
  \item \emph{Recall:} say the Telugu for to inquire, then the Telugu for a refund, then say \emph{saricūsukō} again
\end{itemize}

\begin{tcolorbox}[breakable,colback=teal!4,colframe=teal!35!black,title={Before you move on}]
What does \te{సరిచూసుకో} mean? (\textquotedblleft{}To double-check\textquotedblright{}.) Say it once more.
\end{tcolorbox}
